%%=============================================================================
%% Projectgegevens - VUL DIT ALS EERSTE IN
%%=============================================================================
%%
%% Dit is de enige plek waar je je naam, je titel, je bedrijf en je begeleiders
%% invult. Het voorstel, het rapport en de poster halen hun gegevens hier op,
%% zodat ze nooit uit elkaar kunnen lopen.
%%
%% Enkel de tekst tussen de accolades aanpassen. Laat de commandonamen staan.

%% ---------- Jij ------------------------------------------------------------
\newcommand{\studentnaam}{Zaher Alshikh Kheder}
\newcommand{\studentemail}{zaher.alshikhkheder@student.hogent.be}

%% Wordt gebruikt als bestandsnaam van je PDF's, bijvoorbeeld
%% Familienaam-Voornaam-rapport.pdf. Geen spaties of leestekens gebruiken.
\newcommand{\bestandsnaam}{Alshikhkheder-Zaher}

%% ---------- Je project -----------------------------------------------------
\newcommand{\projecttitel}{Systems monitor}

%% Een ondertitel is optioneel. Laat leeg als je er geen hebt: {}
\newcommand{\projectondertitel}{}

%% Een of twee zinnen. Wordt gebruikt op de poster en in het voorstel.
\newcommand{\projectpitch}{
  Kritieke services van security systemen kan gestopt worden, fowel door een technische storing ofwel dooe een lokale admin op PC, daarom programeer ik een Windows service die die services kan bewaken.NET Core: Ik gebruik de Worker Service-sjabloon samen met het NuGet-pakket Microsoft.Extensions.Hosting.WindowsServices.
}

%% ---------- Waar je het doet ------------------------------------------------
\newcommand{\bedrijfsnaam}{IQGeo}

%% Je mentor op de werkvloer (het bedrijf)
\newcommand{\mentornaam}{Jonas Van Nieuwenberg}
\newcommand{\mentoremail}{jonas.vannieuwenberg@iqgeo.com}

%% Je begeleider van de opleiding (HOGENT)
\newcommand{\begeleidernaam}{Luc Vervoort}
\newcommand{\begeleideremail}{luc.vervoort@hogent.be}

%% ---------- Je repository ---------------------------------------------------
%% De link naar je eigen git-repository. Die hoort zichtbaar in je rapport en
%% op je poster: hij is een beoordeeld onderdeel van je oplevering.
\newcommand{\repolink}{https://github.com/zaheralshikhkheder/GraduaatProject}

%% ---------- Vast voor dit opleidingsonderdeel -------------------------------
%% Deze twee hoef je normaal niet aan te passen.
\newcommand{\opleidingnaam}{Graduaat in het Programmeren}
\newcommand{\academiejaarwaarde}{2026-2027}

%% Zoekwoorden voor het voorstel en de poster (3 tot 5)
\newcommand{\projectkeywords}{trefwoord, trefwoord, trefwoord}
